Equating gravitational force to centrifugal force, we get:
\begin{align*}
\frac{4\pi^2}{T_b^2}\,a_b &= \frac{G\,2M_\odot}{4a_b^2}\\
\frac{T_b^2}{a_b^3} &= 2\,\frac{4\pi^2}{GM_\odot}\\
\Longrightarrow a_b &= 2\ \mathrm{au}.
\end{align*}
\hfill(2.0 pt)

The period of planet can be calculated from the 3rd Kepler's law of motion
\begin{align*}
\frac{T_p^2}{a_p^3} &= \frac{4\pi^2}{GM_\Sigma} = \frac{4\pi^2}{4GM_\odot}\\
\Longrightarrow T_p &= 44.7\ \mathrm{years}
\end{align*}
\hfill(2.0 pt)

Now,
\begin{align*}
\frac{1}{T_s} &= \frac{1}{T_b} - \frac{1}{T_p}\\
\therefore\ T_s &= \frac{T_b T_p}{T_p - T_b} = 4.4\ \mathrm{years}
\end{align*}
\hfill(2.0 pt)

One star always covers exactly the half of the planet's surface. In the best
case, the stars are located at the maximal angular distance from each other.
So the fraction of illuminated surface is simply given by
\[ n = \frac{1}{2} + \frac{\theta}{2\pi}
     = \frac{1}{2} + \frac{2\arctan\!\left(\frac{a_b}{a_p}\right)}{2\pi}
     = 0.53. \]
\hfill(4.0 pt)
